\documentclass[a4paper,12pt]{article}
\usepackage[utf8]{inputenc}
\usepackage[T1]{fontenc}
\usepackage[french]{babel}
\usepackage{geometry}
\geometry{left=2.5cm,right=2.5cm,top=2.5cm,bottom=2.5cm}

\usepackage{lmodern}
\usepackage{xcolor}
\definecolor{titleblue}{RGB}{0,51,140}

\usepackage{hyperref}

\pagestyle{empty}

\begin{document}

\begin{flushleft}
    \textbf{Ismail AISSAOUI} \\
    Ingénieur d'État en Intelligence Artificielle \\
    ismail.aissaoui.pro@gmail.com \\
    +213 660 70 77 96 \\
    \href{https://ismail-aissaoui.vercel.app}{ismail-aissaoui.vercel.app}
\end{flushleft}

\begin{flushright}
    À l'attention de l'équipe d'encadrement du \\
    \textbf{Laboratoire CESI LINEACT} \\
    Campus de Pau \\
    \medskip
    Le 30 Avril 2026
\end{flushright}

\bigskip

\noindent \textbf{Objet : Candidature pour le doctorat en Compression Sémantique pour la XR Aéronautique}

\bigskip

Madame, Monsieur,

Ingénieur d'État en Intelligence Artificielle diplômé de l'ESI-SBA, je vous soumets avec enthousiasme ma candidature pour le doctorat intitulé "Semantic Data Compression for Extended Reality in Aeronautic Industry". Passionné par l'interaction humain-machine et l'optimisation des flux de données pour les environnements immersifs, ce projet soutenu par le Plan France 2030 s'inscrit parfaitement dans ma trajectoire de recherche.

Mon expérience récente m'a permis de développer une expertise technique de pointe sur deux piliers essentiels de cette thèse :

D'une part, mon projet de fin d'études sur le **Jumeau Numérique** pour Sonatrach a été l'occasion de concevoir un système de monitoring industriel intégrant une interface 3D temps réel. J'y ai résolu des problématiques de latence critique en optimisant l'inférence (0.04ms via **C++ JIT xLSTM**) et en concevant une architecture duale Edge/Cloud. Cette compréhension fine des contraintes de performance (Motion-to-Photon delay) est indispensable pour garantir le confort et l'efficacité des opérateurs équipés de casques XR en milieu industriel.

D'autre part, mon travail sur le projet **Geo-RAG** m'a permis de maîtriser les techniques d'extraction sémantique et de raisonnement sur des données multi-sources. Je suis convaincu que l'avenir de la XR industrielle repose sur une compression intelligente "orientée tâche", capable de prioriser les éléments critiques d'une scène 3D. Ma maîtrise de **PyTorch**, de **Unity** et des protocoles de communication distribuée me permettra de concevoir des algorithmes de compression sémantique robustes et économes en énergie.

Intégrer le laboratoire CESI LINEACT à Pau et collaborer au projet Campus Aero Adour (C2A) serait pour moi l'opportunité de porter mes compétences vers un secteur d'excellence mondiale : l'aéronautique. Je suis un candidat autonome, rigoureux et motivé par les défis de la recherche appliquée.

Je vous remercie de l'attention que vous porterez à mon dossier et je me tiens à votre disposition pour un entretien technique.

Je vous prie d'agréer, Madame, Monsieur, l'expression de mes salutations distinguées.

\bigskip

\begin{flushright}
    \textbf{Ismail AISSAOUI}
\end{flushright}

\end{document}
